\subsection{Sprint 4}

Na Sprint 4, já na reta final do projeto, fiquei responsável por desenvolver a tela de aceitar e recusar programas, que serviria como base para a tela de aceitar e recusar instituições. Também auxiliei pontualmente na tela de ajuda e contato. Depois de três sprints trabalhando no \textit{frontend}, eu já estava mais confortável com estilização, mas essa tarefa ainda trouxe novos desafios.

O principal problema técnico apareceu ao usar um \textit{ScrollView} dentro de outro \textit{ScrollView}. No início, o comportamento do scroll interno não funcionava como eu esperava, e precisei pesquisar para entender como resolver. Depois de alguns testes, consegui ajustar o comportamento com uma configuração simples, mas que eu ainda não conhecia. Também implementei um modal de confirmação para que o usuário verificasse se realmente queria aceitar um programa antes de concluir a ação.

Essa sprint foi uma das que mais me fez perceber minha evolução. A tarefa foi mais exigente do que as anteriores, mas eu já tinha repertório para pesquisar, testar e corrigir o problema sem depender tanto dos outros. Na retrospectiva final, ficou claro para mim que o período como AGES I tinha sido muito proveitoso, principalmente por eu ter saído de um conhecimento quase nulo em \textit{frontend} para conseguir entregar telas e componentes reais do projeto.
